\subsection{The actual large-dyadic-conductor Selmer frame}
\label{sec:rh-dyadic-host-frame}

We use the original inertial deformation problem, not the nonresonant
local ring of the good-data argument.  This distinction matters at
the prime $3$.  The following calculation also applies to other odd
coefficient primes at which the finite character is unramified.

Let $N=2^a qd$ be an ordinary RH host conductor, with distinct primes
$q,d\equiv3\pmod4$ greater than $3$.  Let $\psi$ be one of the actual
characters of the host, so $\psi=\psi_2\psi_{\rm odd}$ has order a
power of two, $\psi_{\rm odd}^2=1$, and the conductor of $\psi_2$
is $2^b$.  Fix $e\ge1$ and assume $\psi(-1)=(-1)^e$.
Let $\ell$ be an odd prime not dividing $N$, and suppose
\begin{equation}
                         2^b>8\ell^2.                         \tag{DH1}
\end{equation}
The coefficient ring $O$ is the ring of integers of any finite
extension $P/\mathbf Q_\ell$ containing the values of $\psi$.
Write $k=O/\mathfrak m_O$, $\epsilon$ for the $\ell$-adic cyclotomic
character, $\Psi=\epsilon^e\psi$, and $K=\mathbf Q(i)$ with character
$\eta=\eta_4$.  In the source's notation $\Psi$ is the compensated
character; in particular $\Psi(c)=1$ for complex conjugation $c$.

\begin{theorem}[The original inertial frame and its arithmetic inputs]
\label{thm:rh-dyadic-host-frame}
For the data just specified, construct the original frame $F$ of
\texttt{UnitaryThreeInertia}, using the places over $\ell N$ in
$\mathbf Q$, the quotient $\Gamma=G_{\mathbf Q,S}$ unramified outside
these places, and the characters $\eta,\Psi$ descended to $\Gamma$.
Then the original \texttt{Ok}, \texttt{LocVan}, \texttt{Count},
\texttt{VanE}, and \texttt{Cl1} hold.  The relative dimension defect
in \texttt{Count} is exactly
\[
 d_j=\begin{cases}
       -1&j=-1,\\
       1&j=1,2,\\
       0&\text{otherwise}.
      \end{cases}                                             \tag{DH2}
\]
At $\ell=3$ the actual \texttt{Xpp} holds, for every $e$ of the
stated parity.  Finally, for every nonzero $a_0\in O$, the
$a_0$-torsion of the original strict Selmer group with coefficients
$(P/O)(\Psi^{-1})$ injects into the original $\operatorname{Sha}_O(a_0)$
defined by the inertial frame.  In particular the source's
\texttt{vlen\_le} comparison holds with the same length bound.
No bound for the full integral Selmer length is assumed here.
\end{theorem}

\begin{proof}
\emph{The actual frame and a separating local element.}
Choose a place of the algebraic closure above each prime in $S$;
use its absolute decomposition and inertia groups $D_v,I_v$ and the
canonical map $\pi:G_{\mathbf Q}\to\Gamma$.  Both characters factor
through this quotient and complex conjugation has square one,
$\eta(c)=-1$ and $\Psi(c)=1$.  These are exactly the original
places and character maps.  Since the $\ell$-primary part of $\psi$
is trivial, the source set $\mathrm{Sb}$ of noncoefficient primes
with nontrivial $\ell$-primary inertia is empty.  Thus its set
$\mathrm{Tb}$ is precisely $\{\ell\}$.  Put
\[
        N_v=I_v\cap\ker\eta\cap\ker\overline\Psi .
\]
At $v\ne\ell$ the cyclotomic character is unramified, and
reduction on roots of unity of order prime to $\ell$ is injective.
It follows that $N_v$ is normal of prime-to-$\ell$ index in $I_v$
and $\Psi(N_v)=1$.  These are the exact source subgroups,
including at $v=2$.

For $b\ge3$, the subgroup $1+4\mathbf Z/2^b\mathbf Z$ is cyclic
of order $2^{b-2}$, generated by $5$.  A character of exact conductor
$2^b$ is faithful on that subgroup: its kernel does not contain the
unique subgroup of order two.  For odd $\ell$,
\[
   v_2(\ell^4-1)=v_2(\ell^2-1)+1,\qquad
   2^{v_2(\ell^4-1)}\le2(\ell^2-1)<2^b.
\]
Here the last inequality follows from (DH1).  Consequently
$\psi_2(\ell)^4\ne1$ and, since the odd factor is quadratic,
$\psi(\ell)^4\ne1$.  Choose an arithmetic Frobenius lift in $D_\ell$,
adjusting it by inertia so that its residual cyclotomic value is
one.  This adjustment is possible because
$\mathbf Q_\ell(\mu_\ell)/\mathbf Q_\ell$ is totally ramified with
Galois group $(\mathbf Z/\ell\mathbf Z)^\times$.
Its square $x$ belongs to $\ker\eta$, since $\eta$ has order two.
Unramifiedness of $\psi$ at $\ell$ gives
\begin{equation}
 \bar\epsilon(x)=1,\qquad
 \overline\Psi(x)^2=\overline{\psi(\ell)^4}\ne1.                 \tag{DH3}
\end{equation}
The last inequality again uses the injectivity of reduction on
prime-to-$\ell$ roots of unity.  In particular $\pi(x)$ is an
original marked element.  The inertia character $\eta$ at $\ell$
is trivial since $K$ has discriminant $-4$.
Together with the closedness, normality and continuity of the
local groups, and $D_v/I_v\simeq\widehat{\mathbf Z}$ with its
Frobenius generator and all finite cyclic quotients, these statements
verify every clause of \texttt{Ok}.  The absolute decomposition
groups have not been identified with their images in $\Gamma$.

\emph{The local cohomology and the exact count.}
Let $W_j$ be the $j$-th weight part of the residual adjoint of the
original base representation.  Its matrices have weights given by
differences of $(1,0,-1)$.  On $\ker\eta$ the action is
$\overline\Psi^{\,j}$; on the other coset it is the same scalar
times negative reflection across the antidiagonal.
In particular complex conjugation acts by negative reflection.
Write $n_j=\dim_k W_j$ and $c_j=\dim_k W_j^{c=1}$.  Directly,
\[
\begin{array}{c|rrrrr}
 j&-2&-1&0&1&2\\ \hline
 n_j&1&2&3&2&1\\
 c_j&0&1&1&1&0
\end{array}                                                     \tag{DH4}
\]
and all other spaces vanish.  By (DH3), for $j\ne0$ the element
$x$ has no fixed vector on $W_j$ or on its Tate dual $W_j^*(1)$.
In particular $H^0(D_\ell,W_j)=H^2(D_\ell,W_j)=0$ by local
Tate duality.  The other clause of \texttt{LocVan} is empty.

At any $v\ne\ell$, a local cocycle zero on $N_v$ restricts to a
cocycle of the finite group $I_v/N_v$, since $N_v$ acts trivially.
Averaging in $k$ makes its restriction a coboundary.  Conversely
an unramified class has a representative zero on all inertia;
adding any coboundary preserves vanishing on $N_v$.  Thus the
original local condition on classes at $v\ne\ell$ is the
unramified subspace $L_v=H^1(D_v/I_v,W_j^{I_v})$.
The Frobenius formula gives $\dim L_v=\dim H^0(D_v,W_j)$,
also when inertia acts nontrivially.  At $\ell$ the original
$L_\ell$ is zero for $j<0$, the unramified subspace for $j=0$,
and all of $H^1(D_\ell,W_j)$ for $j>0$.

For clarity, use exactly the source's relative obstruction space:
pairs $(z,(u_v))$ with $z$ a global continuous $2$-cocycle and
$u_v$ a local continuous $1$-cochain with $du_v=\operatorname{res}_v z$,
modulo pairs $(df,(\operatorname{res}_v f))$ and pairs
$(0,(l_v))$ with each $l_v$ an allowed local $1$-cocycle.
Sending a pair to $[z]$ gives the exact sequence
\[
 0\to\operatorname{coker}\left(H^1(\Gamma,W_j)\longrightarrow
                  \bigoplus_v H^1(D_v,W_j)/L_v\right)
 \to\operatorname{Ob}_j\to
 \ker\left(H^2(\Gamma,W_j)\longrightarrow
                  \bigoplus_v H^2(D_v,W_j)\right)\to0.
\]
Indeed a class in the last kernel admits the required local primitives;
a pair with $z=df$ can be reduced to a pair with zero global
component.  Changing $f$ by a global cocycle, or the local
primitives by allowed cocycles, gives exactly the displayed cokernel.
The source Selmer quotient is the kernel of the first displayed map.

All these spaces are finite by the usual finiteness of finite
Galois cohomology.  The global Tate-dual invariants vanish for
$j\ne0$ by (DH3).  For $j=0$ the action is trivial on $G_K$;
since $K\cap\mathbf Q(\mu_\ell)=\mathbf Q$, an element of $G_K$
has nontrivial residual cyclotomic value.  This kills invariants
of $W_0^*(1)$ (also for $\ell=3$).  Poitou--Tate therefore makes
the global $H^2$ restriction onto the direct sum of local
$H^2$ surjective.  Archimedean modified cohomology has no
contribution since $\ell$ is odd.  The exact sequence now gives
\[
 \dim H^1_L-\dim\operatorname{Ob}_j
 =h^1-h^2-\sum_{v\in S}(h_v^1-h_v^2-\dim L_v).                 \tag{DH5}
\]
Here $h^i,h_v^i$ denote dimensions of the indicated global and
local cohomology groups over $k$.
The global Euler characteristic over $\mathbf Q$ is
$h^1-h^2=h^0+n_j-c_j$.
The local one gives $h_v^1-h_v^2=h_v^0$ away from $\ell$
and $h_\ell^1-h_\ell^2=h_\ell^0+n_j$ at $\ell$.
The terms away from $\ell$ cancel.  For $j<0$ the remaining
local term is $n_j$ and the global invariant dimension is zero,
giving $-c_j$.  For $j>0$ the local term is zero (because
$H^2(D_\ell,W_j)=0$), giving $n_j-c_j$.
For $j=0$, global invariants have dimension $c_0$ and the local
term is $n_0$, giving zero.  Substitution of (DH4) gives (DH2),
with the original obstruction quotient, not a different
deformation problem.

\emph{The two global vanishing inputs.}
For \texttt{VanE}, take an original continuous integral cocycle
in weight $j\le0$ with the prescribed local conditions.
After inverting $\ell$, its restrictions away from $\ell$
are unramified: on inertia it is zero on $N_v$, and the remaining
finite quotient has vanishing positive cohomology in characteristic
zero.  At $\ell$, weight zero is unramified by the stated inertia
condition.  The weight-zero representation is trivial over $K$;
there its unramified classes factor through the finite ideal class
group and are zero in a characteristic-zero additive group.
For negative weights restrict to a finite Galois extension
killing $\eta$ and $\psi$.  Each weight becomes a sum of
copies of $\epsilon^{ej}$; the original condition at $\ell$
is a coboundary over the multiplier-one subgroup and hence
over this extension.  Proposition~\ref{prop:selmer-negative-weight-rigidity}
makes every resulting class zero.  In both cases restriction
in characteristic zero is injective by restriction and corestriction.
The original cocycle is consequently an actual coboundary
over the coefficient field.  A nonzero integral scalar clears
the finitely many denominators of its primitive, as required
by \texttt{VanE}.  This argument does not assert that the
integral cohomology is torsion-free.

For \texttt{Cl1}, a locally constant homomorphism from
$\ker(\eta:\Gamma\to\{\pm1\})$ into an arbitrary abelian group
has finite image by compactness.  Pull it back to $G_K$.
At places outside $S$ it is unramified by definition of $\Gamma$;
at the other places the original condition kills every conjugate
of the inertia group intersected with $G_K$.  Hence it factors
through the class group of $K$.  The Gaussian integers are
norm-Euclidean, so this class group is trivial.  The pullback is
zero, and surjectivity onto $\ker\eta$ gives \texttt{Cl1}.

At $\ell=3$, $W_{-2}$ has character
$\bar\gamma=\overline{\eta_4\psi_2^{-2}}$, because the odd factors
and the residual cyclotomic character square to one.
This is an odd character of the full dyadic cyclotomic field
$K_b=\mathbf Q(\mu_{2^b})$.
A cocycle in the original negative-weight local Selmer space
restricts to a homomorphism on $G_{K_b}$.
It vanishes on all inertia away from $3$: on $N_v$ it vanishes,
and the remaining index is prime to $3$.  At $3$ the local
coboundary condition makes its restriction zero on
$G_{K_b,w}$, where the coefficient character is trivial.
The minus $3$-class group vanishes by
Theorem~\ref{thm:dyadic-bernoulli-three-unit} and its class-group
consequence.  The homomorphism is odd under complex conjugation,
so it vanishes on the entire class group.  Restriction on
residual $H^1$ is injective here because $[K_b:\mathbf Q]$
is a power of two.  This proves the literal \texttt{Xpp}.
It works for all positive $e$ in the stated RH frame, not just $e=1$.

\emph{The integral comparison.}
Put $M=(P/O)(\Psi^{-1})$ and $M[a_0]=(a_0^{-1}O/O)(\Psi^{-1})$.
The latter is isomorphic, as a $\Gamma$-module, to
$(O/a_0O)(\Psi^{-1})$.  A strict class on the absolute group
factors through $\Gamma$, since its cocycle vanishes on inertia
outside $S$ where the coefficient action is trivial.
By (DH3), $H^0(\Gamma,M)$ and $H^0(D_\ell,M)$ are both zero:
the action of $x-1$ is multiplication by a unit.
The exact sequence
\[
                    0\to M[a_0]\to M
                         \xrightarrow{a_0}M\to0
\]
therefore gives a unique lift to $H^1(\Gamma,M[a_0])$
of each strict class annihilated by $a_0$.
At $\ell$ its local class is zero by the analogous local
injectivity.  At $v\ne\ell$ the original strict class has local
cocycle a coboundary in $M$; on $N_v$ this is zero because
$\Psi(N_v)=1$.  The lifted cocycle is thus zero there as well.
These are precisely the conditions in the original
$\operatorname{Sha}_O(a_0)$, and the injection follows.
Taking $a_0=\ell^m$ and increasing $m$ exhausts the torsion
group; a uniform bound on the lengths of these subgroups
bounds the length of the union.  This is the original
strict Selmer length, with no loss or division by a residue
degree.
\end{proof}

\begin{theorem}[Crystallinity of the original ordinary points and tangents]
\label{thm:rh-dyadic-host-crystalline}
In the preceding frame, let $r$ be an original characteristic-zero
point for any allowed mark, in the sense of \texttt{IsCharZeroPt}.
Then its three-dimensional matrix representation on $G_K$
is crystalline at the coefficient prime.  Furthermore every
original dual-number lift in \texttt{Rigid} is crystalline on
$G_{K_w}$ as a six-dimensional representation over its coefficient
field, for each $w\mid\ell$.  Its tangent cocycle belongs to the
Bloch--Kato finite local cohomology of the actual polarized adjoint
over $\mathbf Q_\ell$.
\end{theorem}

\begin{proof}
The original point is, up to the permitted torus conjugation,
an integral point reducing to the distinguished representation.
The original \texttt{CondT} provides an ordered full flag on
$G_{K_w}$, after a change of basis congruent to the identity.
Write its diagonal characters as
\[
         \delta_i=\epsilon^{(1-i)e}u_i,\qquad i=0,1,2,
\]
where the $u_i$ are unramified, by $\ell\nmid N$ and the
exact prescribed inertia characters.  They take unit values
in the chosen integral basis.  The element $x$ in (DH3) lies
in $G_{K_w}$ for the chosen local place; the same construction
at each conjugate place gives the corresponding element there.
Since its residual cyclotomic value is one, the three
$\bar u_i(x)$ are exactly the three distinct residual values
$\overline\Psi(x),1,\overline\Psi(x)^{-1}$.
Thus for $i<j$ the character
\[
                      \delta_i\delta_j^{-1}
                    =\epsilon^{(j-i)e}u_i/u_j
\]
has positive Tate exponent, and its Tate dual has no invariants.
For exponent greater than one this follows from infinite
cyclotomic inertia; for exponent one it follows from the
distinct residual unramified values.  The latter is essential
at $\ell=3,e=1$.

An extension of crystalline $V$ by crystalline $U$ is crystalline
if its class in $H^1(G_{K_w},\operatorname{Hom}(V,U))$ belongs
to $H_f^1$.  When the Hom representation has only strictly
negative Hodge--Tate weights (cyclotomic weight $-1$) and its
Tate dual has no invariants, local Euler characteristic and
the Bloch--Kato dimension formula (SP7) give $H^1=H_f^1$.
Apply this first to the first two characters, then to their
extension tensored by $\delta_2^{-1}$.  Its filtration has the
same positive Tate ratios, so the Tate dual still has no
invariants.  This proves the assertion for $r$.

For the dual-number representation the original \texttt{TriIn}
provides an ordered flag by free rank-one dual-number quotients,
with the same exact inertia characters.  Its reduction modulo
the dual-number parameter has the same ordered characters
as the preceding flag.  To see this without imposing a
residual condition on the dual-number change of basis, choose
one inertia element whose three cyclotomic powers are distinct,
possible because $e>0$ and cyclotomic inertia has infinite image.
The invariant subspaces of its diagonalizable matrix with the
specified first one and first two eigenvalues are unique.
Both flags reduce to these subspaces.  Each dual-number
diagonal quotient is an unramified rank-two representation
times the indicated Tate character, since its infinitesimal
character on inertia is trivial.  As an ordinary representation
over the coefficient field it is crystalline and an extension
of $\delta_i$ by itself.  For two such blocks with distinct
indices $i<j$ the Hom representation has a filtration by copies
of $\delta_i/\delta_j$.  The preceding vanishing of Tate-dual
invariants and negativity of Hodge--Tate weights still hold.
The same $H^1=H_f^1$ argument applies to their extensions and
then to the final block.  Hence the entire self-extension is
crystalline.  All representations can be taken over a single
finite coefficient extension: the integral point already is,
and the continuous tangent cocycle takes values in a finite
extension by the cocycle/Baire argument used in
Proposition~\ref{prop:selmer-negative-weight-rigidity} (with
its finite-dimensional adjoint action).
Finally restriction to $G_{K_w}$ is injective on $H^1$ with
coefficients in $B_{\rm cris}\otimes\operatorname{ad}r$,
by restriction and corestriction in characteristic zero.
The crystalline self-extension condition therefore descends
to the actual local adjoint finite condition over $\mathbf Q_\ell$.
\end{proof}

\begin{theorem}[The rigidity and trace suppliers for the inertial case]
\label{thm:rh-dyadic-host-traces}
For the RH data satisfying (DH1), the original implications
\texttt{HostRigid\_s} and \texttt{HostNumber\_s} hold for every
original characteristic-zero point satisfying
\texttt{BPrime\_s} over the Gaussian field $K$.
In particular there is an original element $\phi\in\Gamma$,
independent of the points, and an absolute constant $C$ such that
the trace difference from the distinguished representation is
nonzero at all such points.  Every finite family $(r_i)$ of
such points, pairwise nonisomorphic on $G_K$, satisfies
\[
  \log\ell\sum_i v_\ell\left(
       \operatorname{tr}r_i(\phi)-(\Psi(\phi)+1+\Psi(\phi)^{-1})
                              \right)\le(\ell Ne)^C.           \tag{DH6}
\]
Here $v_\ell(\ell)=1$.  This statement does not supply
\texttt{BPrime\_s} at arbitrary points of the universal ring.
\end{theorem}

\begin{proof}
By its exact definition, \texttt{BPrime\_s} includes irreducibility
on the cyclotomic extension of $K$ and automorphy of
$\rho\epsilon^{-1}$ over $K$, where $\rho=r|_{G_K}$ is the
matrix representation.  More precisely, \texttt{AutOnAlt}
uses this restriction and Tate twist; its
\texttt{IsAutCsdA} condition gives regular algebraic
conjugate-self-dual cuspidal automorphy for the continuous
irreducible representation.  The normalized polarized adjoint
is exactly the adjoint in Thorne's Theorem 6.1, by the same
negative-transpose calculation (TR5) in
Theorem~\ref{thm:support-type-one-rigidity}.  Consequently its
global finite cohomology over $\mathbf Q=K^+$ vanishes.
For an original tangent, the local finite condition at $\ell$
is proved in Theorem~\ref{thm:rh-dyadic-host-crystalline}.
Away from $\ell$, both $r$ and the tangent have the original
condition on $N_v$.  The tangent restricts to a cocycle of
the finite quotient of inertia, since $r(N_v)=1$ and the
cocycle vanishes on $N_v$.  It is a coboundary on inertia
by characteristic-zero averaging.  The tangent class is thus
globally finite, and its vanishing gives the actual coboundary
in \texttt{Rigid}.

For the trace estimate, each such $\rho$ satisfies the original
level condition (SP1), with modulus $N$ and exponents $(e,0,-e)$.
Away from $\ell$, on inertia in $G_K$ the representation
factors through the finite image of $\psi$, because $r$
kills $N_v$.  This abelian finite image is diagonalizable in
characteristic zero and its characters factor through
$\operatorname{Gal}(K(\mu_N)/K)$.  At $\ell$ the exact flag
has the stated exponents and is crystalline, so the first
alternative in (SP1) holds.  The automorphy and irreducibility
requirements of the rank-three class of
Theorem~\ref{thm:support-spectrum} follow from \texttt{BPrime\_s}.
No good-data hypothesis on inertia at $\ell$ has been used.

Choose $p=13$ if $\ell\ne13$ and $p=17$ otherwise.
Both split in $K$; neither is $q$ or $d$, so $p\nmid N$,
and $p\ne\ell$.  Choose $\phi$ to be the image of the inverse
arithmetic Frobenius at one of its places in $K$.
The quotient $\Gamma$ is unramified at $p$, and the prescribed
base trace is
$p^{-e}\psi(p)^{-1}+1+p^e\psi(p)$.
Let $z=\iota(\psi(p)^{-1})$ for the original coefficient
isomorphism $\iota$, and let $\mathcal O$ be the full
Galois orbit of $z$.  Its cardinality is at most $\varphi(N)$.
Apply (SP5)--(SP6) to this orbit, with the given finite family
in the slot for $z$ and empty families elsewhere.  Pairwise
nonisomorphism on $\ker\eta$ is exactly nonisomorphism of
the restrictions $\rho_i$ on $G_K$, because $G_K$ maps
onto $\ker\eta$.  The resulting nonzero integer $B$ gives
\[
 \log\ell\sum_i v_\ell(\operatorname{tr}r_i(\phi)
                         -\operatorname{tr}r_0(\phi))
 \le v_\ell(B)\log\ell\le\log|B|
 \le C_K N[N(e+1)]^{A_K}\log p .
\]
The constants depend only on the fixed $K$ and $p$ is one
of two fixed primes.  Enlarging the exponent gives (DH6)
for every admitted host and every coefficient ring.  Applying
the nonvanishing conclusion of (SP6) to a singleton gives
the stated nonvanishing for every point with the same $\phi$.
\end{proof}

\paragraph{Exact remaining integral frontier.}
At $\ell=3$ the original large-conductor exponent-one case
has $2^b\ge9216$, hence (DH1).  Theorems
\ref{thm:rh-dyadic-host-frame}--\ref{thm:rh-dyadic-host-traces}
supply the original local and global arithmetic inputs to
\texttt{abstract\_count}, including its strict Selmer lower
comparison and the uniform trace upper bound.
They also apply to high-dyadic-conductor characters for
$e\ge2$ at $3$.  If the original \texttt{ThmS6\_s}
(finiteness of the weight-zero part of the actual universal
ring and automorphy of its type-I points) and the original
graded algebra estimate in \texttt{AbstractCount} are
established for this inertial problem, its output is
\[
     V_{\rm len}\le V\le2e_O\sum_i
           v_\ell(\operatorname{tr}r_i(\phi)-\operatorname{tr}r_0(\phi)).
\]
Dividing by the original ramification index $e_O$ and
multiplying by $\log\ell$ gives the polynomial bound for
the actual normalized Selmer length, using (DH6).
This last inference is conditional on those integral
universal-ring and graded-algebra inputs; the earlier
good-data ring cannot simply be substituted for this one.
For $\ell\ne3$ the residual weight-$-2$ condition remains
to be established when the source requires it.  For general
Dirichlet GRH the nontrivial odd-conductor factors must also
be retained.  The weighted product and the full W04
obligation are not closed by the present calculation.
